% Example "heat": the problem, the decay of the modes in space-time, the
% stability limit against the grid, and the time and space errors. Every
% number comes from data/, written by examples/heat/main.f90: run
% `fpm run --example heat` from the project root, then build this file with
% lualatex in examples/heat/.
\documentclass[11pt,a4paper]{article}
\usepackage[margin=2.2cm]{geometry}
\usepackage{amsmath,booktabs,array}
\usepackage{pgfplots}
\usepgfplotslibrary{groupplots,colormaps}
\pgfplotsset{compat=1.18}
\pagestyle{empty}

\colorlet{m-euler}{black}
\colorlet{m-order2}{orange!85!black}
\colorlet{m-order3}{green!50!black}
\colorlet{m-rk4}{blue!80!black}
\colorlet{m-dp54}{red!75!black}
% Five viridis colours, one per time in the profiles.
\definecolor{vir1}{HTML}{440154}
\definecolor{vir2}{HTML}{3B528B}
\definecolor{vir3}{HTML}{21918C}
\definecolor{vir4}{HTML}{5EC962}
\definecolor{vir5}{HTML}{C8B100}
% Diverging colour map for u: blue below zero, white at zero, red above.
\pgfplotsset{colormap={heat}{rgb255(0cm)=(33,102,172) rgb255(1cm)=(247,247,247)
  rgb255(3cm)=(244,165,130) rgb255(5.4cm)=(103,0,13)}}

\begin{document}
\section*{Example: the heat equation, method of lines}

\subsection*{1. Problem}
\begin{equation*}
  u_t = u_{xx}, \quad 0 < x < 1, \qquad u(0, t) = u(1, t) = 0, \qquad
  u(x, 0) = \sin\pi x + \tfrac12\sin 4\pi x + \tfrac14\sin 9\pi x .
\end{equation*}
Each sine mode decays on its own: $u = \sum_k a_k e^{-(k\pi)^2 t}\sin k\pi x$.
Central differences on $n$ interior nodes, $\Delta x = 1/(n + 1)$, turn it
into $u' = Au$, $A$ the tridiagonal matrix $(1, -2, 1)/\Delta x^2$, whose
eigenvectors are the same sine modes with eigenvalues
\begin{equation*}
  \lambda_k = -\frac{4}{\Delta x^2}\sin^2\frac{k\pi\Delta x}{2},
  \qquad -\frac{4}{\Delta x^2} < \lambda_n \le \lambda_k \le \lambda_1 \approx -\pi^2 ,
\end{equation*}
so the semi-discrete system is solved exactly by the same sum with
$\lambda_k$ in place of $-(k\pi)^2$. The fastest mode $\lambda_n$ grows like
$\Delta x^{-2}$: the system becomes stiff as the grid is refined. The example
uses the model classes: the state holds the boundary nodes too, which the
model's \texttt{apply\_boundaries} sets to zero before every evaluation.
\begin{center}
\small
% Written by examples/heat/main.f90: do not edit.
\begin{tabular}{l>{\raggedright\arraybackslash}p{12.4cm}}
\toprule
initial condition & $u(x, 0) = \sin\pi x + \tfrac12\sin 4\pi x + \tfrac14\sin 9\pi x$ \\
space & central differences, $n$ interior nodes, $\Delta x = 1/(n + 1)$; boundary nodes set to 0 by \texttt{apply\_boundaries} \\
time convergence & $n = 63$, error against the semi-discrete solution at $t = 0.02$, $\Delta t = 10^{-4}, \dots, 10^{-4}/64$; order from the two finest levels above $10^{-13}$, kept if within 0.3 of the next slope \\
stability & $n = 15, 31, 63, 127, 255$; largest $\Delta t$ for which 3000 steps do not let $\max|u|$ grow (initial value plus $10^{-6}(-1)^j$); bisection in $[0.8, 1.2]$ times the prediction; table: the ratio farthest from 1 \\
PDE error & at $t = 0.02$ with $\Delta t \approx 0.9$ of each method's limit; order from the two finest grids \\
figures & \texttt{rk4}, $n = 63$ \\
methods & \texttt{rk3ws}, \texttt{rk4}, \texttt{dp54} \\
interface & model classes, \texttt{rk\_scheme}, \texttt{apply\_boundaries} for the boundary \\
\bottomrule
\end{tabular}

\end{center}

\newpage
\subsection*{2. Solutions}
\begin{center}
\begin{tikzpicture}
  % ---- (a) space-time ----
  \begin{axis}[
      name=map,                                            % anchor for the panels on the right
      width=8cm, height=10.6cm,                            % tall map
      view={0}{90}, xmin=0, xmax=1, ymin=0, ymax=0.01,     % x across, t up: the fast modes fade here
      xlabel={$x$}, ylabel={$t$}, ytick={0,0.002,...,0.01}, % space-time
      scaled y ticks=false, yticklabel style={/pgf/number format/fixed, /pgf/number format/precision=3}, % 0.002, not 10^-3
      title={(a) $u(x, t)$, \texttt{rk4}, $n = 63$},       % panel title
      colormap name=heat, point meta min=-0.25, point meta max=1.35, % the range of u
      colorbar horizontal, colorbar style={at={(0.5,-0.13)}, anchor=north, xlabel={$u$}, height=0.2cm, width=6cm},
      enlargelimits=false, axis on top,                    % the map fills the axes
    ]
    \addplot3[surf, shader=interp, mesh/ordering=x varies] table[x=x, y=t, z=u] {data/spacetime.dat};
  \end{axis}
  % ---- (b) profiles ----
  \begin{axis}[
      name=profiles,                                       % anchor for (c)
      at={($(map.north east)+(1.5cm,0cm)$)}, anchor=north west, % right of (a)
      width=6.8cm, height=6.2cm,                           % panel
      xmin=0, xmax=1, ymin=-0.35, ymax=1.45,               % the profiles
      xlabel={$x$}, ylabel={$u$},                          % solution against space
      title={(b) profiles: \texttt{rk4} (dots), exact PDE (lines)}, % panel title
      axis background/.style={fill=black!2},               % very light background
      legend style={font=\tiny, at={(0.98,0.98)}, anchor=north east, fill=white, fill opacity=0.85, text opacity=1, draw=black!30},
    ]
    \addplot[vir1, thick, no marks] table[x=x, y=e1] {data/profiles.dat};            % exact PDE solution, t = 0
    \addlegendentry{$t = 0$}
    \addplot[vir1, only marks, mark=*, mark size=0.9pt, forget plot]
      table[x=x, y=u1] {data/profiles.dat};                                    % rk4, n = 63, t = 0
    \addplot[vir2, thick, no marks] table[x=x, y=e2] {data/profiles.dat};            % exact PDE solution, t = 0.002
    \addlegendentry{$t = 0.002$}
    \addplot[vir2, only marks, mark=*, mark size=0.9pt, forget plot]
      table[x=x, y=u2] {data/profiles.dat};                                    % rk4, n = 63, t = 0.002
    \addplot[vir3, thick, no marks] table[x=x, y=e3] {data/profiles.dat};            % exact PDE solution, t = 0.005
    \addlegendentry{$t = 0.005$}
    \addplot[vir3, only marks, mark=*, mark size=0.9pt, forget plot]
      table[x=x, y=u3] {data/profiles.dat};                                    % rk4, n = 63, t = 0.005
    \addplot[vir4, thick, no marks] table[x=x, y=e4] {data/profiles.dat};            % exact PDE solution, t = 0.01
    \addlegendentry{$t = 0.01$}
    \addplot[vir4, only marks, mark=*, mark size=0.9pt, forget plot]
      table[x=x, y=u4] {data/profiles.dat};                                    % rk4, n = 63, t = 0.01
    \addplot[vir5, thick, no marks] table[x=x, y=e5] {data/profiles.dat};            % exact PDE solution, t = 0.03
    \addlegendentry{$t = 0.03$}
    \addplot[vir5, only marks, mark=*, mark size=0.9pt, forget plot]
      table[x=x, y=u5] {data/profiles.dat};                                    % rk4, n = 63, t = 0.03
  \end{axis}
  % ---- (c) stability against the grid ----
  \begin{axis}[
      at={($(profiles.south west)+(0,-1.6cm)$)}, anchor=north west, % below (b)
      width=6.8cm, height=6.2cm,                           % panel
      xmode=log, ymode=log, xlabel={$\Delta x$}, ylabel={largest stable $\Delta t$}, % log-log
      xtick={0.004,0.01,0.03,0.06}, xticklabels={0.004,0.01,0.03,0.06},             % plain grid spacings
      title={(c) stability limit, measured and predicted}, % panel title
      axis background/.style={fill=black!2},               % very light background
      legend style={font=\tiny, at={(0.02,0.98)}, anchor=north west, fill=white, fill opacity=0.85, text opacity=1, draw=black!30},
    ]
    \addplot[m-rk4, thin, forget plot] table[x=dx, y=predicted] {data/stable_rk4.dat};                  % h* = 2.79
    \addplot[m-rk4, only marks, mark=diamond*, mark size=1.8pt] table[x=dx, y=measured] {data/stable_rk4.dat};
    \addlegendentry{\texttt{rk4}}
    \addplot[m-dp54, thin, forget plot] table[x=dx, y=predicted] {data/stable_dp54.dat};                % h* = 3.31
    \addplot[m-dp54, only marks, mark=pentagon*, mark size=1.8pt] table[x=dx, y=measured] {data/stable_dp54.dat};
    \addlegendentry{\texttt{dp54}}
  \end{axis}
\end{tikzpicture}
\end{center}
(a) The first hundredth of a time unit: the nine bumps of $\sin 9\pi x$ fade
first (decay rate $81\pi^2$, gone by $t \approx 0.004$), then the four of
$\sin 4\pi x$ ($16\pi^2$), leaving the slow $\sin\pi x$. (b) Profiles at five times: \texttt{rk4} on 63 nodes (dots) on the
exact PDE solution (lines). (c) The largest stable step on five grids
(markers) against the prediction $h^*/|\lambda_n| \approx h^*\Delta x^2/4$
(lines, slope 2): every method follows it within $0.25\,\%$; methods with
the same real-axis limit share a line.

\newpage
\subsection*{3. Tests: time error, stability, space error}

\paragraph{Prediction.} (i) Against the exact semi-discrete solution only the
time error remains: each method converges at its order on this linear
problem (\texttt{rk3ws}: 3). (ii) A step is stable while $h|\lambda_n| \le
h^*$, so $\Delta t_{\max} = h^*/|\lambda_n| \approx h^*\Delta x^2/4$.
(iii) Against the PDE, with $\Delta t$ near that limit, the time error is of
order $\Delta t^p \sim \Delta x^{2p}$, far below the space error
$O(\Delta x^2)$: every method should give the same second-order error in
$\Delta x$.

\paragraph{Result.}
\begin{center}
\small
% Written by examples/heat/main.f90: do not edit.
\begin{tabular}{lccccc}
\toprule
& & \multicolumn{2}{c}{time convergence, $n = 63$} & stability & PDE \\
\cmidrule(lr){3-4}
name & $h^*$ & order & observed & measured / predicted & order in $\Delta x$ \\
\midrule
\texttt{rk3ws} & 2.5127 & 2 & 2.99 & 1.0011 &  2.00 \\
\texttt{rk4} & 2.7853 & 4 & 4.01 & 1.0011 &  2.00 \\
\texttt{dp54} & 3.3066 & 5 & --- & 1.0008 &  2.00 \\
\bottomrule
\end{tabular}

\end{center}

\begin{center}
\begin{tikzpicture}
  \begin{groupplot}[
      group style={group size=2 by 1, horizontal sep=1.8cm}, % two panels
      width=7.4cm, height=6.6cm,                           % size of each panel
      xmode=log, ymode=log, grid=major, grid style={black!8}, % log-log
      axis background/.style={fill=black!2},               % very light background
      every axis plot/.append style={thick, mark size=1.3pt}, % same weight for all curves
    ]
    % ---- (i) time convergence ----
    \nextgroupplot[
      title={(i) time error, $n = 63$, $t = 0.02$}, xlabel={$\Delta t$}, % against the semi-discrete solution
      legend to name=htlegend, legend columns=7,           % legend below the panels
      legend style={font=\small, /tikz/every even column/.append style={column sep=0.25cm}},
    ]
    \addplot[m-order3, dashed, mark=triangle] table[x=dt, y=error] {data/error_rk3ws.dat};
    \addlegendentry{\texttt{rk3ws}}
    \addplot[m-rk4, mark=diamond*]            table[x=dt, y=error] {data/error_rk4.dat};
    \addlegendentry{\texttt{rk4}}
    \addplot[m-dp54, mark=pentagon*]          table[x=dt, y=error] {data/error_dp54.dat};
    \addlegendentry{\texttt{dp54}}
    % ---- (ii) PDE error ----
    \nextgroupplot[title={(ii) error against the PDE, $t = 0.02$}, xlabel={$\Delta x$}]
    \addplot[m-order3, dashed, mark=triangle] table[x=dx, y=error] {data/pde_rk3ws.dat};
    \addplot[m-rk4, mark=diamond*]            table[x=dx, y=error] {data/pde_rk4.dat};
    \addplot[m-dp54, mark=pentagon*]          table[x=dx, y=error] {data/pde_dp54.dat};
    \addplot[black, thin, densely dotted, domain=0.0035:0.07] {1.0*x^2};          % slope 2 guide
    \node[font=\scriptsize, anchor=north west] at (axis cs:0.012,1.4e-4) {$\Delta x^2$};
  \end{groupplot}
  \node[below=0.2cm] at (current bounding box.south)        % legend centred under the panels
    {\pgfplotslegendfromname{htlegend}};
\end{tikzpicture}
\end{center}
(i) Every method converges at its order (\texttt{rk3ws}: 3 on this linear
problem). \texttt{dp54} is at round-off already at the largest step tested,
$10^{-4}$. (ii) Against the PDE the three methods give the same error, of
order $\Delta x^2$ ($1.5\times10^{-5}$ on the finest grid for all three): at
the steps stability allows, the time error is negligible. For a stiff, diffusive problem the explicit step is set
by stability, not by accuracy, and a higher order buys almost nothing.
\end{document}
